\section{Saturation and the differential-inductance matrix}
An analytic local illustration makes the nonlinear content of the symmetry
visible:
\begin{equation}
 \coenergy=W_0(\id)+a_2(\id)\iq^2+a_4(\id)\iq^4+\cdots.
 \label{eq:local-energy}
\end{equation}
This expansion illustrates parity; $C^2$ regularity alone does not guarantee
an analytic series, and the series is not our general machine model. Its
derivatives are
\begin{align}
 \psid&=W_0'(\id)+a_2'(\id)\iq^2+a_4'(\id)\iq^4+\cdots,\\
 \psiq&=2a_2(\id)\iq+4a_4(\id)\iq^3+\cdots.
\end{align}
PM bias and d saturation enter $W_0$; q saturation changes the dependence on
$i_q$; cross-saturation appears when the q coefficients depend on $i_d$.
All can coexist without producing a forbidden parity component. Energy-based
nonlinear models have also been used to estimate saturation experimentally
\cite{jebai2011}. Our argument here is a structural differentiation of
\cref{eq:energy-parity}, not a fitted constant-inductance approximation.

For reciprocity we need to know how flux changes when its current
argument moves. A single secant inductance cannot describe that displacement.
The appropriate local linear map is the Jacobian
\begin{equation}
 \Ldiff=\nabla^2\coenergy=\frac{\partial\flux}{\partial\current}
 =\begin{bmatrix}L_{dd}&L_{dq}\\L_{qd}&L_{qq}\end{bmatrix},
 \label{eq:jacobian}
\end{equation}
where
\begin{align}
 L_{dd}&=\partial_{i_d}\psid,&L_{dq}&=\partial_{i_q}\psid,\\
 L_{qd}&=\partial_{i_d}\psiq,&L_{qq}&=\partial_{i_q}\psiq.
\end{align}
It maps an infinitesimal current displacement to a flux displacement. Since
it is the Hessian of a $C^2$ co-energy, equality of mixed derivatives gives
\begin{equation}
 \boxed{L_{dq}=L_{qd}.}\label{eq:reciprocity}
\end{equation}
This reciprocity is about differential cross-inductances, not equality of
$L_{dd}$ and $L_{qq}$. Another view is an infinitesimal closed loop in current
space: a conservative flux gradient has zero circulation, so the two mixed
changes must cancel. Independently fitted flux surfaces may violate this
condition even if each surface looks smooth.

Geometrically, the gradient of co-energy gives flux and its Hessian gives
curvature. A quadratic linear-machine surface has constant curvature.
Saturation changes that curvature with current; cross-saturation changes
the mixed curvature, coupling the two directions. Differential inductance
is therefore a matrix of local curvatures rather than two globally fixed
scalar inductances.

Differentiating an even function with respect to $i_d$ retains its parity;
differentiating it with respect to $i_q$ reverses it. For an odd function the
roles reverse. Hence
\begin{equation}
 L_{dd},L_{qq}\ \text{are even in }i_q,\qquad
 L_{dq},L_{qd}\ \text{are odd in }i_q.\label{eq:inductance-parities}
\end{equation}
In particular the cross terms vanish on $i_q=0$ for a smooth symmetric map.
Nonlinear differential gains remain permitted everywhere else. Reciprocity
requires conservative magnetics and sufficient smoothness; parity requires
the additional reflection assumption. Neither should be imposed on a
dynamic loss model merely because it is convenient.
